\section{Introduction}
As LLM training shifted to internet-scale data, stereotype learning moved from an edge case to a predictable byproduct of scale \citep{nadeem2021stereoset,nangia2020crowspairs,blodgett2021salmon}. The dominant practical response has been post-training local editing \citep{vig2020causal,zou2023repeng,turner2023steering,tan2024steering}, but paired causal write-versus-erase tests and transfer stress tests are still less common than benchmark-local evaluations \citep{meade2022survey,morabito2023debiasing,jin2021transferability,reusens2023multilingual}.

We ask a direct causal question: after stereotype-related behavior is learned, are linear residual-direction local edits enough to reliably remove it? Our audit combines matched inject-versus-remove interventions, component-level diagnostics, and cross-source transfer stress tests in one protocol. The central result is asymmetric controllability: writing stereotype preference is easier than erasing it. In the confirmatory core, completion-format cross-position matched contrast is significant in 2/3 models, while a same-position-valid autoregressive setting is significant in 3/3. In exploratory cross-position extension models, support is heterogeneous (1/4 score-significant; 3/4 margin-significant). Results establish cross-position causal leverage of trait-position directions at decision time, not same-position representational erasure at the decision token.
This paper does not compare mitigation efficacy across methods; it audits the inferential validity of local linear-direction removal claims.

Two practical boundaries are important up front. First, completion-format tests and same-position autoregressive tests probe related but distinct causal objects, so we report both rather than collapsing them into a single effect estimate. Second, model-level estimates are pair-weighted over a source-skewed included slice, and strongest headline counts depend on mixed-source direction construction; they should not be interpreted as source-invariant guarantees. These constraints do not erase the asymmetry result, but they bound where the result should be applied.

\section{Related Work}
Prior work established benchmark families and measurement caveats for stereotype bias (including StereoSet and CrowS-Pairs) \citep{nadeem2021stereoset,nangia2020crowspairs,blodgett2021salmon}, and showed that debiasing gains can vary by benchmark and objective \citep{meade2022survey,morabito2023debiasing}. Mechanistic editing work shows that many behaviors are writable through activation steering and linear directions \citep{vig2020causal,zou2023repeng,turner2023steering}, but this does not by itself establish that the same local objects are reliably erasable.

Attribution and causal-tracing methods support component ranking, but sign quality and intervention directionality vary across settings \citep{kramar2024atp,syed2024atp,elhage2021mathematical}. Transfer studies also show uneven carryover across tasks and domains \citep{jin2021transferability,reusens2023multilingual}. Our contribution is a causal-and-transfer audit that directly tests writeability, eraseability, and robustness under matched operators.

\section{Setup and Evaluation Protocol}
This protocol asks whether detected stereotype directions are causally writable and removable at decision time. Primary evaluations use StereoSet intrasentence completion pairs and transfer stress tests use CrowS-Pairs. Core runs use deterministic 480/120 train/test splits per model; exact source-axis counts and alignment coverage are reported in Appendix Tables~\ref{tab:counts}, \ref{tab:alignment-retention-main}, and \ref{tab:coverage-bias}.

Coverage is intentionally defined with two denominators. ``Pre-sampling alignability'' counts how many raw benchmark pairs satisfy single-token alignment before any held-out cap is applied. ``Post-sampling inclusion'' counts how many of those alignable pairs enter the frozen analysis slice used for core comparisons. We report both because they answer different questions: alignability bounds format validity, while inclusion determines the actual evidential mix behind each model-level estimate. In the frozen Rev3 slice, source composition is skewed relative to raw pools, so aggregated model-level claims should be read as pair-weighted summaries of the included slice rather than source-balanced estimates of all raw benchmark content.

For pair $i$ at prediction position $p_i$, we score
\[
m_i = \log p(\text{stereo}_i \mid x_i, p_i) - \log p(\text{anti}_i \mid x_i, p_i),
\]
with behavioral endpoint $\mathbf{1}[m_i>0]$ and mechanistic endpoint $m_i$. In completion format, stereotype/anti contexts are identical up to $p_i$, so prediction-position stereotype directions are zero by construction; we therefore extract directions at trait positions and intervene at prediction position (Appendix Table~\ref{tab:dirpos-portability}). Same-position AR replication is analyzed separately as a complementary setting.

This distinction matters for interpretation. In completion format, cross-position intervention can still have real causal leverage at the decision token, but it is not evidence of same-position representational erasure at that token. In same-position AR prompts, the edited signal can be present natively at decision time, so ``remove'' is a stricter removability test. Reporting both settings prevents a common inferential shortcut where a single format is treated as universally representative of all deployment-time editing conditions.

\paragraph{Causal logic of the audit.}
Many local debiasing claims implicitly use a three-step chain: identify a stereotype-linked internal signal, edit it, then infer that behavior is robustly debiased. We test that third step explicitly. Injection is a sufficiency/writeability probe (does adding direction $d$ increase stereotype preference?), while ablation is a necessity/eraseability probe (does removing $d$ reliably reduce it?). If injection is strong and ablation weak, the correct interpretation is one-sided controllability, not successful local removal.

\paragraph{Method operators.}
At layer $\ell$ and axis $a$, direction ablation projects out $d_{\ell,a}$ at prediction position:
\(h_{\ell,p} \leftarrow h_{\ell,p} - (h_{\ell,p}\cdot \hat d_{\ell,a})\hat d_{\ell,a}\).
Matched-direction injection adds $+d_{\ell,a}$ at prediction position in anti-stereotype runs; we use it as a diagnostic positive control, not a mitigation recommendation. For ranked-component diagnostics, we also evaluate single-site and multi-site ablations over DLA/AtP candidate sets.

\paragraph{Direction construction and portability.}
Directions are defined by train-split difference-in-means at trait-token positions:
\[
d_{\ell,a} = \mu^{\text{stereo}}_{\ell,a} - \mu^{\text{anti}}_{\ell,a}.
\]
In completion format, extracting at prediction position is degenerate (zero vectors by construction), so usable directions must be trait-position objects applied cross-position at decision time (Appendix Table~\ref{tab:dirpos-portability}). This design constraint is why we separately report same-position AR replication as a distinct causal setting rather than assuming direct equivalence.

Main inferential claims use paired tests with explicit correction-family labeling. The preregistered family was Bonferroni; the manuscript primary rows report within-table BH values from frozen Exp16 outputs for consistency with the run-native tables, and we additionally report across-model BH plus preregistered Bonferroni robustness in Appendix Table~\ref{tab:inference-robustness}. Core directional conclusions are unchanged across these families. We treat preregistered same-position assumptions that are invalid under completion-format constraints as non-adjudicable and document this in the preregistration map (Appendix Table~\ref{tab:prereg-map}).

All runs are deterministic, inference-only (no training/fine-tuning), and reported with fixed artifact pointers. Core and extension models span 2B--7B open-weight decoders. The frozen Rev3 run family used for manuscript artifacts was executed on Quadro RTX 8000 48GB and Quadro RTX 5000 16GB GPUs, with approximately 5.1 total GPU-hours across the seven model lanes.

\paragraph{Protocol correction and provenance.}
The preregistered draft assumed same-position direction extraction in the primary completion-format setting. That assumption is not valid for StereoSet-style intrasentence completion because stereotype/anti contexts are identical through the scored token. We therefore treat completion-format cross-position analysis as the confirmatory family for that format, and same-position AR as a separate, complementary family with explicit interaction checks (Appendix Table~\ref{tab:cross-vs-same-interaction}). This is a methodological correction, not an outcome-driven endpoint swap, and the preregistered-to-final mapping is explicit in Appendix Table~\ref{tab:prereg-map}.

We also separate computational provenance from inferential definition. Table~\ref{tab:matched-contrast} is unchanged in metric definition (paired inject-minus-remove score contrast), but the frozen values come from the Rev3 run map rather than the older 2026-05-08 family. The old-versus-frozen mapping and correction-family recomputation are reported in Appendix Tables~\ref{tab:table2-provenance} and \ref{tab:inference-robustness}.

\section{Main Causal Result: Writing Is Easier Than Erasing}
\label{sec:main-result}

This is the core causal test in the paper: if local directions are truly removable bottlenecks, removal should be comparably effective to injection under matched operators.

\paragraph{Cross-position design.}
In completion format, direction extraction and intervention operate at different positions: directions are estimated at trait positions $t_i$, but interventions are applied at prediction position $p_i$. Prediction-position directions are zero by construction in this format (Appendix Table~\ref{tab:dirpos-portability}), so Exp16 is a cross-position stress test of causal leverage/removability rather than a same-position erasure test. We therefore report same-position AR results (Exp26) as a complementary setting rather than a direct format replication.

Table~\ref{tab:main-intervention-score} gives the compact intervention outcomes: injection score deltas are significant in Gemma and Llama and positive-but-non-significant in Gemma-IT, while ablation is mixed (negative in all three models, significant only in Llama). The canonical asymmetry endpoint is the matched contrast in Table~\ref{tab:matched-contrast}, computed from paired per-item inject-minus-remove deltas.

\begin{table}[t]
\centering
\footnotesize
\setlength{\tabcolsep}{2.5pt}
\begin{tabular}{@{}l p{0.21\columnwidth} c p{0.21\columnwidth} c@{}}
\toprule
Model & Ablation $\Delta$score [95\% CI] & $q$ & Injection $\Delta$score [95\% CI] & $q$ \\
\midrule
Gemma-2-2B & $-$0.058 [$-$0.108, $-$0.008] & 0.0654 & +0.075 [0.033, 0.125] & 0.00781 \\
Gemma-2-2B-IT & $-$0.008 [$-$0.067, 0.042] & 1.000 & +0.042 [0.008, 0.075] & 0.125 \\
Llama-3.2-3B & $-$0.100 [$-$0.158, $-$0.050] & 0.000488 & +0.233 [0.133, 0.342] & 0.0000509 \\
\bottomrule
\end{tabular}
\caption{Main intervention score outcomes ($n=120$ pairs/model). Full baseline$\rightarrow$post rows are reported in Appendix Table~\ref{tab:main-intervention-score-full}. This table uses full-sequence residual injection and prediction-position ablation; the clean matched write-vs-erase test is Table~\ref{tab:matched-contrast}.}
\label{tab:main-intervention-score}
\end{table}

\begin{table}[tbp]
\centering
\footnotesize
\setlength{\tabcolsep}{3pt}
\begin{tabular}{@{}l p{0.48\columnwidth} c@{}}
\toprule
Model & Inject $-$ Remove $\Delta$score [95\% CI] & $q$ \\
\midrule
Gemma-2-2B & +0.133 [0.075, 0.192] & 0.000145 \\
Gemma-2-2B-IT & +0.050 [0.000, 0.100] & 0.146 \\
Llama-3.2-3B & +0.333 [0.242, 0.433] & 0.00000002 \\
\bottomrule
\end{tabular}
\caption{Canonical matched-operator asymmetry contrast from the matched-operator contrast test, using paired per-item contrasts $c_i=\delta_i^{\text{inject}}-\delta_i^{\text{remove}}$. $q$ values are within-table BH-FDR from the frozen Exp16 outputs; across-model BH and preregistered Bonferroni robustness are reported in Appendix Table~\ref{tab:inference-robustness}.}
\label{tab:matched-contrast}
\end{table}

\begin{figure}[t]
\centering
\includegraphics[width=\columnwidth]{build/fig_main_asymmetry_rev3.png}
\caption{Core asymmetry view from the Rev3 frozen values. Injection is consistently positive, ablation is weak, and inject-minus-remove contrast is largest for Llama and Gemma.}
\label{fig:main-asymmetry}
\end{figure}

Because baseline stereotype score is below chance in all three core models in this frozen slice, we interpret injection as increasing stereotype preference rather than assuming strong baseline stereotyping. The largest mechanistic movement in the core set remains Llama injection margin shift (+1.278 in Appendix Table~\ref{tab:main-intervention-score-margin}), while score-level removability remains limited. This pattern motivates a write-versus-erase interpretation rather than ratio-style claims.

To reduce base-condition artifacts, we also inspect matched 2$\times$2 anti/stereo-base rows (Appendix Table~\ref{tab:matched-matrix}). These rows are consistent with Table~\ref{tab:matched-contrast}: remove-on-anti is small in Gemma and Gemma-IT but larger for Llama, while inject-on-anti is consistently positive and strongest in Llama. We therefore treat the paired inject-minus-remove contrast as the primary behavioral adjudication.

The paired contrast is the main inference target because it compares inject/remove outcomes on the same items. Under the table-reported within-run BH family, Gemma and Llama are significant; Gemma-IT is directionally aligned but non-significant ($q=0.146$). The preregistered Bonferroni family yields the same directional core verdict (Gemma/Llama significant, Gemma-IT non-significant; Appendix Table~\ref{tab:inference-robustness}). We keep score as the primary endpoint and margin as sensitivity evidence because injection often moves both score and margin, while ablation can move margin without many decision flips.

\paragraph{Claim A (cross-position completion stress test).}
Exp16 is a format-constrained cross-position stress test in which directions are trait-position objects applied at prediction time. In the core set, matched score asymmetry is confirmed in 2/3 models (Gemma and Llama); in exploratory extension runs under the same cross-position family, support is heterogeneous (1/4 score-significant; 3/4 margin-significant; Appendix Table~\ref{tab:ext4-summary}). Because model-level estimates are pair-weighted over a source-skewed included slice, these are mixed-slice operational outcomes rather than source-invariant population estimates.

\paragraph{Claim B (same-position AR removability test).}
Exp26 is a distinct same-position-valid autoregressive setting on overlapping pair IDs with stereotype-relevant context moved into the prefix before the scored token. Here, matched score asymmetry is significant in 3/3 core models (Appendix Table~\ref{tab:exp26-same-position}), giving stronger removability stress evidence than the completion-format condition.

\paragraph{Bridge test (interaction, not pooled estimand).}
Exp16 and Exp26 probe related but non-identical causal objects, so we do not pool them into one aggregate fraction. Instead, we report the formal same-minus-cross interaction (Appendix Table~\ref{tab:cross-vs-same-interaction}): Llama is positive and significant, Gemma is slightly negative, and Gemma-IT is positive but non-significant.

\paragraph{Extension instability interpretation.}
Extension-model asymmetry is exploratory and should not be read as confirmatory generalization. The Rev3 family shows attenuation/reversal in some added models (notably OLMo), and several extension rows are low-$n$ ($n=44$ or $80$), so variance from run-family lineage and slice composition is non-trivial. We therefore use extension rows as stress evidence for heterogeneity boundaries, not as a standalone external-validity proof.

\paragraph{Direction occupancy and removability.}
To separate ``zero mean direction'' from ``zero per-example occupancy,'' we report occupancy diagnostics over pair-level hidden states in both settings. In completion-format cross-position tests, prediction-position occupancy of trait-derived directions is centered near zero at the aggregate level by construction, but individual examples still have non-zero projections. In same-position AR, occupancy magnitude is typically larger and more behavior-coupled. The occupancy tables (Appendix Tables~\ref{tab:occupancy-setting} and \ref{tab:occupancy-coupling}) quantify this shift and show that larger baseline $|h\cdot \hat d|$ tends to predict larger ablation margin movement, with random-direction baselines used to bound pure norm effects.

\paragraph{Geometry without flips.}
Ablation often moves margin geometry without reliably changing binary score decisions. Boundary-crossing diagnostics (Appendix Tables~\ref{tab:boundary-crossing}, \ref{tab:boundary-crossing-quality}) show that many pairs are far from the decision boundary, so moderate confidence shifts can coexist with few behavioral flips. This is one reason margin-level sensitivity can exceed score-level sensitivity and why we keep score as the primary endpoint for confirmatory interpretation.

Across all three core models, boundary crossings under ablation concentrate in near-boundary items (negative $\rho(|m_0|,\mathbf{1}\{\text{cross}\})$), while far-from-boundary items show confidence movement without flips. This is why margin-level asymmetry is easier to satisfy than score-level asymmetry in the cross-position run family: margin responds to directional push, score requires boundary crossing.

\paragraph{Specificity and composition constraints.}
Injection-specificity controls show that true-direction injection usually exceeds norm-random controls and often exceeds shuffled-axis controls (Appendix Table~\ref{tab:injection-controls}), reducing a pure-norm perturbation explanation. However, composition-sensitivity reruns show attenuation under source-specific direction pools (Appendix Tables~\ref{tab:composition-sensitivity}, \ref{tab:composition-sensitivity-interaction}), so headline asymmetry support should be read as conditional on mixed-source direction construction.

Composition effects are material, not cosmetic. For Gemma, the mixed-pool contrast (+0.133) attenuates to +0.048 and non-significance under StereoSet-only direction training, while Llama shows a significant CrowS-only attenuation relative to mixed training (Appendix Table~\ref{tab:composition-sensitivity-interaction}). We therefore avoid framing mixed-pool outcomes as source-invariant stereotype semantics. The practical reading is narrower: mixed-pool direction construction can yield stronger mixed-slice operational asymmetry in this protocol, but source-specific direction pools can weaken the effect.
Axis-reweighted and leave-one-axis-out companions show bounded but non-zero aggregate sensitivity (Appendix Tables~\ref{tab:axis-reweighted}, \ref{tab:axis-loao}), reinforcing that pooled claims should be read as slice-conditioned rather than fully axis-invariant.

Span-level robustness is similarly bounded. Exp27 (ablation-only) is a sensitivity check, not a full matched asymmetry test; Exp28 provides matched span-level inject-minus-remove contrasts. For Gemma and Gemma-IT, Exp28 span rows are CrowS-dominated by available aligned pairs, so those rows should not be generalized as broad cross-source span evidence. Llama has wider source coverage in span scoring and retains a significant matched contrast, but this should be read as partial support rather than universal span-level confirmation (Appendix Tables~\ref{tab:exp27-multitoken}, \ref{tab:exp28-multitoken-matched}, \ref{tab:exp28-source-strata}).

Finally, interpretation as stereotype-semantic causality is conditional on plausibility/fluency confounds not dominating benchmark labels. Our controls constrain generic perturbation explanations and axis-agnostic random-direction explanations, but they do not fully isolate stereotype semantics from possible benchmark-internal plausibility structure.

Prompt-based calibration rows in the appendix provide an additional operator-contrast boundary: prompt matched contrasts are substantially larger than direction matched contrasts in the same endpoint family (Appendix Tables~\ref{tab:prompt-matched-contrast}, \ref{tab:prompt-vs-direction-ratio}, \ref{tab:matched-contrast}). This does not invalidate the direction-edit audit; it narrows the claim. Our asymmetry evidence is strongest for linear residual-direction projection/editing, not for all local mitigation operators, and should be interpreted as method-family-specific.

\FloatBarrier
\section{Why Erasing Is Brittle}
\label{sec:brittle}

Attribution-sign diagnostics and multi-site stress tests indicate that erasure is brittle even when localization signal is detectable. In the legacy diagnostic family, Gemma-IT shows below-chance DLA sign agreement (13/46, $q=0.0136$), while Gemma and Llama are near chance; under strict split-clean reruns, that legacy Gemma-IT inversion reverses to chance (31/61, $q=1.000$), so we treat it as a reproducibility-risk diagnostic rather than a standalone confirmatory claim (Appendix Table~\ref{tab:split-clean-diagnostics}). Multi-site suppressor contamination is non-zero across models (about 13.8\%--26.3\%), implying that sign-blind top-$k$ sets often include components that push opposite to intended debiasing direction (Appendix Tables~\ref{tab:split-clean-diagnostics}, \ref{tab:multisite-heterogeneity}).

Axis-level diagnostics reinforce a heterogeneity interpretation rather than a global sign verdict. Gemma race DLA is an aligned boundary-condition case (5/5 agreement, $q=0.004$), while many other axis-model cells are near chance or non-evaluable at current component counts (Appendix Table~\ref{tab:sign-audit-axis}). This supports a practical rule: attribution sign can occasionally be usable in narrow axis/model slices, but model-level editing pipelines should not assume sign reliability without direct causal adjudication.

Rank-sweep diagnostics are non-monotonic as lower-ranked components are added, consistent with suppressor contamination/basis misalignment rather than a clean single-subspace removal story (Appendix Tables~\ref{tab:multisite-full}, \ref{tab:alt-mech}). Practically, ranking methods remain useful candidate generators, but sign-driven erasure should still be paired with direct causal checks.

\FloatBarrier
\section{Transfer Stress Test}
\label{sec:transfer}

\textbf{Scope note.} Cross-source per-axis sample sizes are small ($n=2$--13), so this section is powered to reject only large transfer effects.

\begin{figure*}[t]
\centering
\includegraphics[width=0.9\textwidth]{build/fig_transfer_ablation_rev3.png}
\caption{Core ablation transfer heatmap for the condition-level matrix. The dominant pattern is weak/near-zero cross-source behavior with one within-source negative trend in Gemma CrowS$\rightarrow$CrowS.}
\label{fig:transfer-heatmap}
\end{figure*}

\begin{table*}[t]
\centering
\footnotesize
\setlength{\tabcolsep}{3pt}
\begin{tabular}{@{}lllccl@{}}
\toprule
Model & Type & Condition & $n_{\text{pairs}}$ & Ablation $\Delta$score [95\% CI] & $q$ \\
\midrule
Gemma-2-2B & Within & StereoSet$\rightarrow$StereoSet & 39 & +0.051 [$-$0.128, 0.256] & 1.000 \\
Gemma-2-2B & Cross & StereoSet$\rightarrow$CrowS & 13 & +0.077 [$-$0.154, 0.308] & 1.000 \\
Gemma-2-2B & Within & CrowS$\rightarrow$CrowS & 77 & $-$0.143 [$-$0.247, $-$0.039] & 0.0768 \\
Gemma-2-2B & Cross & CrowS$\rightarrow$StereoSet & 19 & +0.000 [$-$0.211, 0.211] & 1.000 \\
\midrule
Gemma-2-2B-IT & Within & StereoSet$\rightarrow$StereoSet & 39 & +0.077 [$-$0.077, 0.231] & 0.625 \\
Gemma-2-2B-IT & Cross & StereoSet$\rightarrow$CrowS & 13 & +0.385 [0.154, 0.692] & 0.250 \\
Gemma-2-2B-IT & Within & CrowS$\rightarrow$CrowS & 77 & $-$0.039 [$-$0.130, 0.052] & 0.625 \\
Gemma-2-2B-IT & Cross & CrowS$\rightarrow$StereoSet & 19 & +0.105 [$-$0.053, 0.316] & 0.625 \\
\midrule
Llama-3.2-3B & Within & StereoSet$\rightarrow$StereoSet & 46 & +0.022 [$-$0.109, 0.152] & 1.000 \\
Llama-3.2-3B & Cross & StereoSet$\rightarrow$CrowS & 11 & +0.091 [$-$0.182, 0.364] & 1.000 \\
Llama-3.2-3B & Within & CrowS$\rightarrow$CrowS & 70 & +0.086 [$-$0.071, 0.243] & 1.000 \\
Llama-3.2-3B & Cross & CrowS$\rightarrow$StereoSet & 26 & +0.000 [$-$0.231, 0.192] & 1.000 \\
\bottomrule
\end{tabular}
\caption{Condition-level ablation transfer matrix (main rows; full decompositions in Appendix Table~\ref{tab:transfer-matrix}).}
\label{tab:transfer-main}
\end{table*}

The ablation transfer matrix is mostly null at current power: no cross-source row survives FDR, and even the strongest within-source trend (Gemma CrowS$\rightarrow$CrowS: $\Delta$score $=-0.143$, $q=0.0768$) is non-significant (Appendix Tables~\ref{tab:transfer-matrix}, \ref{tab:transfer-axis}, \ref{tab:transfer-equivalence}). This means Step 4 in this draft is primarily a stress test for null stability and backfire risk, not a clean ``successful-within-source then failed-cross-source'' transfer adjudication.

The strongest transfer boundary condition is not a null but a reversal: in positive-control injection transfer, Llama StereoSet$\rightarrow$CrowS is significantly negative ($-0.350$, $q=0.0208$; Appendix Table~\ref{tab:exp30-transfer}). This implies source-conditioned direction geometry can invert across benchmarks. A plausible mechanism is benchmark-specific direction misalignment: source-conditioned stereotype directions show substantial negative cosine/sign alignment, so writing in one benchmark geometry can move anti-directionally in another (Appendix Tables~\ref{tab:exp30-backfire-axis}, \ref{tab:exp30-direction-alignment}).

Rows with $q=\mathrm{n/a}$ reflect zero-variance/test-degenerate paired deltas, not missing files. Overall interpretation is conservative: within-source ablation is mostly weak at current power, but positive-control injection shows that benchmark shift can still produce backfire, so Step 4 remains a practical risk screen.

\section{Synthesis and Practical Validation Standard}
\label{sec:synthesis}

Across settings, linear residual-direction edits are frequently writable without being cleanly erasable. Claim A (cross-position completion stress test) is confirmed in 2/3 core models, with heterogeneous exploratory extension support (1/4 score-significant, 3/4 margin-significant). Claim B (same-position AR removability test) is confirmed in 3/3 core models. We do not pool A and B into a single estimand. All model-level summaries are pair-weighted over source-skewed included slices.

The results are strongest as a method-specific claim: for linear direction projection/editing, control leverage does not imply reliable eraseability. They do not establish asymmetry for all mitigation mechanisms, and they are not source-invariant given composition sensitivity and transfer fragility. In this draft, Step 1 is strong while Steps 2--4 are partly provisional, so we frame the checklist as a falsification protocol (minimum audit standard), not a certification pass.

For deployment-facing local-edit claims, we recommend a minimum four-step protocol:
\begin{enumerate}[leftmargin=*]
  \item \textbf{Causal efficacy check}: direct inject-minus-remove contrast under matched operators.
  \item \textbf{Reliability screen}: attribution-sign audit before sign-driven editing.
  \item \textbf{Contamination screen}: suppressor-rate audit in candidate intervention sets.
  \item \textbf{Transfer check}: cross-benchmark/format validation before generalization claims.
\end{enumerate}
Passing this protocol is necessary but not sufficient; in this draft, not all steps are strongly passed simultaneously (Appendix Table~\ref{tab:checklist-status}). Operationally: if Step 1 passes but Steps 2--4 fail or remain provisional, treat the intervention as control-only, not as a debiasing guarantee.

\section{Conclusion}
This audit answers the central question for linear residual-direction edits: stereotype preference is more reliably writable than erasable under matched interventions. The key implication is inferential: finding a stereotype-linked direction does not by itself justify a removal claim.

Two caveats shape the final interpretation. First, completion-format cross-position asymmetry and same-position AR asymmetry are complementary, not interchangeable, and model-by-setting interactions show they need separate reporting. Second, extension and mixed-pool outcomes are conditional on source composition and should not be over-read as source-invariant behavior. Within those bounds, the practical message is stable: local editing is a useful control tool, but not a debiasing guarantee. In practice, when Step 1 passes but Steps 2--4 do not, the intervention should be used only as controllability tooling, not as validated debiasing.

\section{Limitations}
\begin{enumerate}[leftmargin=*]
  \item \textbf{Model scope and scale.} The audit covers seven open-weight decoder models (roughly 2B--7B). Results may differ in larger or closed-weight systems, including stronger RLHF/DPO deployment regimes.
  \item \textbf{Benchmark and format scope.} Evaluation is completion-style, which enforces cross-position intervention at the decision token; same-position conclusions in alternative formats are not established.
  \item \textbf{Transfer power limits.} Cross-source axis-level rows are low-$n$, so current transfer tests can rule out only unusually large effects and cannot resolve moderate transfer magnitudes.
  \item \textbf{Mechanism and method scope.} Pathway diagnostics are coarse (non-identifying at circuit resolution), and the study evaluates linear residual-direction editing rather than head-to-head mitigation efficacy across methods.
\end{enumerate}

\section{Ethical Considerations}
This paper studies stereotype-related behavior and mitigation diagnostics in existing language models. Interventions are used for causal auditing and validation, not as a deployment recommendation by themselves.

All datasets used here (StereoSet, CrowS-Pairs, SEEGeL) and all model weights evaluated (Gemma-2, Llama-3.2, Qwen-3, Mistral-7B, OLMo-2) are publicly released open research artifacts used under their respective licenses for research use. This paper reports derived statistics and does not redistribute third-party dataset or model artifacts. AI assistants were used for code-completion and editorial polishing of draft text; experimental design, analysis decisions, and scientific claims were made and verified by the authors.
